% Vietnamese translation for OI Public Library
% Source-File: IOI中国国家候选队论文/2008/Day1/2.郑暾《平衡规划——浅析一类平衡思想的应用》/平衡规划.ppt
\documentclass[11pt]{article}
\usepackage{oipl-vn}

\newcommand{\Oh}{\mathcal{O}}

\title{Bản trình chiếu: Quy hoạch cân bằng}
\author{Zheng Dun}
\date{2008}

\begin{document}
\maketitle

\origtitle{Slide companion for balanced planning}
\sourcepath{IOI China candidate team papers / 2008 / Day 1 / Zheng Dun.ppt}

\section{Phạm vi bản dịch}

Tệp PowerPoint đi cùng bài viết đã được dịch từ tài liệu Word. Bản ghi chú này
gom lại tư tưởng cân bằng trong thuật toán, cấu trúc dữ liệu và hiện thực
chương trình.

\section{Điểm cân bằng}

Trong phòng thi, lựa chọn thuật toán không chỉ dựa vào độ phức tạp lý thuyết.
Một thuật toán rất mạnh nhưng khó cài có thể kém thực dụng hơn một biến thể
đơn giản hơn, đặc biệt khi dữ liệu không chạm tới trường hợp xấu nhất. Tư
tưởng cân bằng là tìm điểm giữa tốc độ, bộ nhớ, độ khó lập trình và độ chắc
chắn của lời giải.

\section{Các mẫu thường gặp}

Bài viết nêu ba mẫu: hiện thực không điển hình của thuật toán kinh điển, dùng
thuật toán không hoàn hảo nhưng đủ tốt, và cấu trúc hóa đơn giản cho một bài
toán phức tạp. Ví dụ về ghép cặp tổng quát cho thấy đôi khi ta tránh một thuật
toán khó như blossom bằng một chiến lược ngẫu nhiên hoặc thực nghiệm có kiểm
soát, nếu mục tiêu là vượt qua dữ liệu thi cụ thể.

\section{Kết luận}

Quy hoạch cân bằng không thay thế chứng minh thuật toán khi cần lời giải hoàn
chỉnh. Nó là một kỹ năng lựa chọn: biết khi nào nên bỏ thêm bộ nhớ để giảm
thời gian, khi nào nên dùng thuật toán đơn giản hơn, và khi nào nên chuyển
phần khó sang một bước xử lý nhỏ hơn nhưng chắc chắn đúng.

\end{document}
